\begin{abstract}
\textbf{Objective:}
Covariance fidelity, source-trained predictive likelihood, and
discriminative transfer optimize different regularisation targets in
cross-subject EEG. We examine whether objective disagreement changes
unknown-person classification and whether more exact source
optimisation or source-contrast robustness resolves the discrepancy.
\textbf{Methods:}
We study source-fitted tangent-space regularised LDA under a
zero-target-information protocol, correct a validation-class-centred
covariance score to source-mean Gaussian predictive likelihood,
and examine exact pairwise rank transitions, conditional
spectral shrinkage bounds, and an exact decomposition of the
proper predictive mean-shift penalty into AUC-invariant
participant location and class-contrast components. An exploratory source-only
contrast-uncertainty-envelope criterion was evaluated against
the original grouped 27-point source-AUC selector on an
unchanged classifier path. All target results were retained.
\textbf{Results:}
In a frozen EEGMMIDB block ($N=20$), discriminative source
selection exceeded the historical class-centred score by
0.0561 AUC (Holm $p=0.00597$), but five people lost more
than 0.02 AUC. A protocol-first Weibo2014 study ($N=10$)
showed no target difference between continuous and discrete
source optimisation (both macro AUC 0.73141).
In Cho2017, source choices were frozen using 12 source people
before examining target people 13--32: the contrast-envelope
criterion improved over the source-AUC grid by only
0.003395 AUC ($N=20$, exact sign-flip $p=0.0931$).
Without retuning, a subsequent, explicitly sequential
20-person block reversed the difference to $-0.003500$;
the combined 40-person difference was $-0.000053$ AUC
(descriptive bootstrap 95\% interval
$[-0.003218,0.002973]$). The envelope selected
full shrinkage $\lambda=1$, algebraically identical to a fixed
full-shrinkage LDA classifier.
\textbf{Conclusion:}
The prediction--classification objective mismatch includes
a nonnegative participant-location term to which the
within-person rank AUC is exactly invariant. More precise or
conditionally robust source selection did not produce stable
transfer gains in these tests. Explicit
validation estimands, strict selection boundaries and
honest adverse-person evidence are essential; no universally
superior regularisation rule is established.
\end{abstract}
